\chapter{Goldbach}\label{conjectures:ch:goldbach}

\begin{center}\itshape
That every even number is a sum of two primes, I hold to be a completely certain theorem, although I cannot prove it.
\upshape(Leonhard Euler to Christian Goldbach, 30 June 1742)
\end{center}

On 7 June 1742 Christian Goldbach, a German scholar in the Russian service, wrote a long letter to his friend Leonhard Euler in Berlin. Most of
it is about series. Near the end Goldbach risks a guess about sums of primes, and in the margin, as if the idea came to him while rereading, he adds a
sharper one: ``it seems at least that every number greater than $2$ is a sum of three primes''. Goldbach counted $1$ as a prime, as many did then. Euler
answered on 30 June. He replied that the guess would follow from a cleaner statement which Goldbach had once told him, that every even number is a sum of
two primes, and he gave the verdict in the epigraph: certain, but unproved.

It is still unproved. Every even number up to $4\cdot10^{18}$ has been checked by computer; the conjecture is part of Hilbert's eighth problem; in 2000
the publishers Faber and Faber offered a million dollars for a proof within two years, to promote a novel, and kept their money. This chapter explains
both halves of that story: why everybody believes the conjecture, and why nobody can prove it. On the way we prove the first theorem that said
anything at all about every integer: Lev Schnirelmann's theorem of 1930 that every integer greater than $1$ is a sum of a bounded number of primes.
The proof combines a sieve upper bound from Chapter~\ref{conjectures:ch:twins} with a disarmingly simple idea about densities.

\begin{think}
How many ways are there to write $100$ as a sum of two primes? And $98$? Which do you expect to have more representations, $2^{20}$ or $3\cdot2^{19}$?
\end{think}

\section{The comet}

For $n\ge1$ let
\[
r(n)=\#\{(p,q):\ p,q\text{ primes},\ p+q=n\},
\]
counting ordered pairs, so $r(100)=12$ from the six unordered sums $3+97$, $11+89$, $17+83$, $29+71$, $41+59$, $47+53$. For odd $n$ one of the primes
must be $2$, so $r(n)$ is $0$ or $2$ and odd $n$ are uninteresting (Exercise~\ref{conjectures:ex:oddr}). Goldbach's conjecture, in Euler's form, says
$r(n)>0$ for every even $n\ge4$.

Plot $r(n)$ against $n$ and you see Figure~\ref{conjectures:fig:comet}, nicknamed the \emph{Goldbach comet}. The values do not merely stay positive: they grow, roughly
like $\frac n{(\log n)^2}$, and the smallest of them pulls away from $0$ faster and faster. At $n=10^6$ there are $10\,804$ ordered representations.
The comet has a structure too: it is split into bands, and the even numbers divisible by $3$ have about twice as many representations as the others.

\begin{figure}[ht]
\centering
\fig{conjectures-ch03-comet}
\caption{The Goldbach comet: $r(n)$ for $2500$ randomly chosen even $n\le10^6$. Blue: $3\nmid n$; green: $3\mid n$ but $15\nmid n$; red: $15\mid n$. The
dashed curves are the Hardy--Littlewood prediction for $n$ a power of $2$, and twice and $\frac83$ times it, the lower edges of the three bands.}
\label{conjectures:fig:comet}
\end{figure}

The bands are explained by the same heuristic as the twin prime constants of Chapter~\ref{conjectures:ch:twins}. Pick $k<n$ at random. Naively, $k$ and $n-k$ are
prime with probability about $\frac1{\log k}\cdot\frac1{\log(n-k)}$. Now correct for each prime $p$. If $p\nmid n$, the two numbers $k$ and $n-k$ avoid
the multiples of $p$ with probability $1-\frac2p$, because $k$ must avoid the two classes $0$ and $n$; independent numbers would do so with probability
$(1-\frac1p)^2$. If $p\mid n$, the two classes coincide, and the probability is $1-\frac1p$. The correction factor is the product of the ratios,
\[
\mathfrak S(n)=\prod_{p\mid n}\frac{1-\frac1p}{(1-\frac1p)^2}\prod_{p\nmid n}\frac{1-\frac2p}{(1-\frac1p)^2}=2C_2\prod_{\substack{p\mid n\\p>2}}\frac{p-1}{p-2}\qquad(n\text{ even}),
\]
where $C_2=\prod_{p>2}\bigl(1-\frac1{(p-1)^2}\bigr)=0.6601618\ldots$ is the twin prime constant: for odd $p\nmid n$ the ratio is $1-\frac1{(p-1)^2}$, for odd
$p\mid n$ it is $\frac p{p-1}=\bigl(1-\frac1{(p-1)^2}\bigr)\frac{p-1}{p-2}$, and for $p=2$ it is $2$.

\begin{remark}[Hardy and Littlewood, 1923]\label{conjectures:rem:hlgoldbach}
Hardy and Littlewood conjectured that, for even $n\to\infty$,
\[
r(n)\sim\mathfrak S(n)\,J(n),\qquad J(n)=\int_2^{n-2}\frac{\dd t}{\log t\,\log(n-t)}\sim\frac n{(\log n)^2}.
\]
This is a conjecture, not a theorem: nobody can prove even $r(n)>0$. Our derivation above is a heuristic, and we never use it in a proof.
\end{remark}

The factor $\prod_{p\mid n,p>2}\frac{p-1}{p-2}$ is $2$ when $3\mid n$, and $2\cdot\frac43$ when $15\mid n$: those are the bands (Exercise~\ref{conjectures:ex:bands}).
Dividing $r(n)$ by the prediction collapses the comet to a thin line at height $1$ (Figure~\ref{conjectures:fig:normalised}). At $n=10^6$, with
$\mathfrak S(10^6)=2C_2\cdot\frac43=1.7604$, the prediction is $10\,834$ against $10\,804$ observed, a ratio of $0.997$.

\begin{figure}[ht]
\centering
\fig{conjectures-ch03-normalised}
\caption{The comet collapses: $r(n)/\bigl(\mathfrak S(n)J(n)\bigr)$ for the same $2500$ even $n$, colours as in Figure~\ref{conjectures:fig:comet}. Beyond $10^5$ every
value lies between $0.95$ and $1.05$; the four smallest $n$ of the sample fall outside the window.}
\label{conjectures:fig:normalised}
\end{figure}

\paragraph{How far it has been checked.} Computer checks reached $4\cdot10^{11}$ in 1993 (Matti Sinisalo) and $4\cdot10^{14}$ in 2001 (J\"org Richstein). The current record is due to Tom\'as
Oliveira e Silva, Siegfried Herzog and Silvio Pardi, whose distributed computation reached $4\cdot10^{18}$ on 4 April 2012 (double-checked to
$4\cdot10^{17}$; published in 2014). For each even $n$ they found the least prime $p(n)$ with $n-p(n)$ prime. It is always tiny: the largest value they
met was $p(n)=9781$, for $n=3\,325\,581\,707\,333\,960\,528$. Below $10^6$ the record is $p(503\,222)=523$. Finding a Goldbach partition is never the
hard part; the hard part is knowing that one exists.

\section{Schnirelmann's density}

In 1930 Lev Schnirelmann, a 25-year-old Soviet mathematician, found a way to say something about every integer at once. His idea is to stop
asking for two primes and ask for a bounded number of them, and to replace the question ``is $n$ represented?'' by a question about densities that
behaves well under addition.

For a set $A$ of nonnegative integers, write $A(n)=\#\{a\in A:\ 1\le a\le n\}$, counting from $1$. The \emph{Schnirelmann density} of $A$ is
\[
\sigma(A)=\inf_{n\ge1}\frac{A(n)}n.
\]
It is a harsh density. If $1\notin A$ then $\sigma(A)=0$, whatever the rest of $A$ is, because $\frac{A(1)}1=0$. And $\sigma(A)=1$ exactly when $A$ contains every
positive integer (Exercise~\ref{conjectures:ex:sdensity}). The harshness is the point: a lower bound $A(n)\ge\sigma(A)\,n$ holds for every $n$, not just eventually.
For two sets, $A+B=\{a+b:\ a\in A,\ b\in B\}$, and $hA=A+\dots+A$ ($h$ terms).

\begin{theorem}[Schnirelmann's inequality]\label{conjectures:thm:schnirlemma}
Let $A$ and $B$ be sets of nonnegative integers containing $0$, with $\sigma(A)=\alpha$ and $\sigma(B)=\beta$. Then
\[
\sigma(A+B)\ge\alpha+\beta-\alpha\beta,\qquad\text{that is,}\qquad1-\sigma(A+B)\le(1-\alpha)(1-\beta).
\]
\end{theorem}
\begin{proof}
Fix $n\ge1$, and let $0=a_0<a_1<\dots<a_k\le n$ be the elements of $A$ in $[0,n]$, so $k=A(n)$. Put $a_{k+1}=n+1$. Between consecutive elements lie the
gaps $g_i=a_{i+1}-a_i-1\ge0$ for $0\le i\le k$, and $\sum_{i=0}^kg_i=n-k$, since the gaps together contain every integer of $[1,n]$ not in $A$.

Every $a_i$ with $i\ge1$ lies in $A+B$, because $0\in B$. Moreover, for each $i$ and each $b\in B$ with $1\le b\le g_i$, the number $a_i+b$ lies in
$A+B$ and strictly between $a_i$ and $a_{i+1}$, so it is in the $i$-th gap. These numbers are distinct for different $(i,b)$, and distinct from the
$a_i$. Hence
\[
(A+B)(n)\ge k+\sum_{i=0}^kB(g_i)\ge k+\beta\sum_{i=0}^kg_i=k+\beta(n-k)=(1-\beta)k+\beta n.
\]
Since $k\ge\alpha n$ and $1-\beta\ge0$, this is at least $(1-\beta)\alpha n+\beta n=(\alpha+\beta-\alpha\beta)n$. Dividing by $n$ and taking the infimum
gives the theorem.
\end{proof}

\begin{lemma}\label{conjectures:lem:cover}
Let $A$ and $B$ contain $0$, with $\sigma(A)+\sigma(B)\ge1$. Then every nonnegative integer lies in $A+B$.
\end{lemma}
\begin{proof}
Let $n\ge1$ (and $0=0+0$). If $n\in A$ or $n\in B$ we are done, since the other set contains $0$. Otherwise $A(n-1)=A(n)\ge\sigma(A)n$ and likewise for
$B$, so $A(n-1)+B(n-1)\ge n$. The sets $\{a\in A:1\le a\le n-1\}$ and $\{n-b:\ b\in B,\ 1\le b\le n-1\}$ both lie in $[1,n-1]$, which has $n-1$
elements, and their sizes add up to at least $n$. So they meet: $a=n-b$ for some $a\in A$, $b\in B$.
\end{proof}

\begin{corollary}\label{conjectures:cor:basis}
If $0\in A$ and $\sigma(A)>0$, there is $h$ such that every nonnegative integer is a sum of $h$ elements of $A$.
\end{corollary}
\begin{proof}
By Theorem~\ref{conjectures:thm:schnirlemma} and induction, $1-\sigma(kA)\le(1-\sigma(A))^k$ for every $k\ge1$. Choose $k$ with $(1-\sigma(A))^k\le\frac12$. Then
$\sigma(kA)\ge\frac12$, and Lemma~\ref{conjectures:lem:cover} applied to $kA+kA$ shows that $h=2k$ works.
\end{proof}

The corollary turns a question about every integer into a question about a single number, the density. A set with positive Schnirelmann density is
an additive basis of finite order. Schnirelmann's plan for the primes was to show that the sums of two primes, which are too sparse to touch every
even number as far as anyone can prove, are at least dense enough to have positive density.

\section{Every integer is a sum of boundedly many primes}

We need an upper bound for $r(m)$, and it is the bound of Proposition~\ref{conjectures:prop:pairs} in disguise.

\begin{proposition}\label{conjectures:prop:goldbachsieve}
There is an absolute constant $C$ such that for every even $m\ge4$,
\[
r(m)\le C\Bigl(\frac m{\varphi(m)}\Bigr)^2\frac m{(\log m)^2}.
\]
\end{proposition}
\begin{proof}
The number $r(m)$ counts the primes $p\le m$ with $m-p$ prime. Sieve the set $\mathcal A=\{k(m-k):\ 1\le k\le m\}$ by all primes up to $z$, as in the
proof of Proposition~\ref{conjectures:prop:pairs}. A prime $p>z$ with $m-p>z$ prime gives an element of $\mathcal A$ with no prime factor up to $z$, so
$r(m)\le2z+S(\mathcal A,\mathcal P,z)$. For squarefree $d$, the number of classes $k\bmod d$ with $d\mid k(m-k)$ is multiplicative, equal to $1$ at
primes $p\mid m$ (the classes $0$ and $m$ coincide) and $2$ at the other primes; the error in counting the $k\le m$ in these classes is at most their
number. These are exactly the data of the proof of Proposition~\ref{conjectures:prop:pairs}, with $h$ replaced by $m$ (even) and $x$ by $m$, and that proof used
nothing else. It gives the bound.
\end{proof}

The factor $\bigl(\frac m{\varphi(m)}\bigr)^2$ is unbounded, but it is bounded on average, even raised to a higher power.

\begin{lemma}\label{conjectures:lem:phifour}
There is a constant $C_3$ such that $\displaystyle\sum_{m\le n}\Bigl(\frac m{\varphi(m)}\Bigr)^4\le C_3\,n$ for every $n\ge1$.
\end{lemma}
\begin{proof}
Write $\bigl(\frac m{\varphi(m)}\bigr)^4=\prod_{p\mid m}(1+f(p))=\sum_{d\mid m}\mu^2(d)f(d)$, where $f(p)=\bigl(1+\frac1{p-1}\bigr)^4-1$, extended
multiplicatively to squarefree $d$. With $u=\frac1{p-1}\le1$, $f(p)=u(4+6u+4u^2+u^3)\le15u\le\frac{30}p$. Exchanging the sums,
\[
\sum_{m\le n}\Bigl(\frac m{\varphi(m)}\Bigr)^4=\sum_{d\le n}\mu^2(d)f(d)\Bigl\lfloor\frac nd\Bigr\rfloor\le n\sum_d\frac{\mu^2(d)f(d)}d=n\prod_p\Bigl(1+\frac{f(p)}p\Bigr)\le n\prod_p\Bigl(1+\frac{30}{p^2}\Bigr),
\]
and the last product converges.
\end{proof}

\begin{theorem}[Schnirelmann, 1930]\label{conjectures:thm:schnirelmann}
There is an integer $S$ such that every integer $n\ge2$ is a sum of at most $S$ primes.
\end{theorem}
\begin{proof}
Let $A=\{0,1\}\cup\{p+q:\ p,q\text{ primes}\}$. We show $\sigma(A)>0$.

\emph{Counting with multiplicity.} Let $n$ be large. On the one hand, the sum $R(n)=\sum_{m\le n}r(m)$ counts the pairs of primes with $p+q\le n$, which
include all pairs with $p,q\le\frac n2$. By the prime number theorem (Tome~5, Theorem~5.5), $\pi(\frac n2)\ge\frac n{3\log n}$ for large $n$, so
\[
R(n)\ge\pi\bigl(\tfrac n2\bigr)^2\ge\frac{n^2}{9(\log n)^2}.
\]
On the other hand, the function $t\mapsto\frac t{(\log t)^2}$ increases for $t\ge e^2$, so Proposition~\ref{conjectures:prop:goldbachsieve} gives
$r(m)\le C\bigl(\frac m{\varphi(m)}\bigr)^2\frac n{(\log n)^2}$ for even $m$ with $8\le m\le n$. For odd $m$, $r(m)\le2$, and $r(m)\le2$ also for even $m<8$.
With Lemma~\ref{conjectures:lem:phifour},
\[
\sum_{m\le n}r(m)^2\le C^2\frac{n^2}{(\log n)^4}\sum_{m\le n}\Bigl(\frac m{\varphi(m)}\Bigr)^4+4n\le C^2C_3\frac{n^3}{(\log n)^4}+4n\le C_4\frac{n^3}{(\log n)^4}.
\]

\emph{Cauchy--Schwarz.} Let $T(n)$ be the number of $m\le n$ with $r(m)>0$. Then $R(n)=\sum_{m\le n,\,r(m)>0}r(m)\le\sqrt{T(n)}\bigl(\sum_{m\le n}r(m)^2\bigr)^{1/2}$,
so
\[
T(n)\ge\frac{R(n)^2}{\sum_{m\le n}r(m)^2}\ge\frac{n^4/(81(\log n)^4)}{C_4n^3/(\log n)^4}=\frac n{81C_4}.
\]
Every such $m$ lies in $A$, so $A(n)\ge\frac n{81C_4}$ for all $n\ge n_0$. For $1\le n<n_0$, $A(n)\ge1\ge\frac n{n_0}$ because $1\in A$. Hence
$\sigma(A)\ge\min\bigl(\frac1{81C_4},\frac1{n_0}\bigr)>0$.

\emph{From $A$ to primes.} By Corollary~\ref{conjectures:cor:basis}, there is $h$ such that every nonnegative integer is a sum of $h$ elements of $A$. Let $n\ge2$ and write
$n-2=a_1+\dots+a_h$ with $a_i\in A$. Each $a_i$ is $0$, $1$ or a sum of two primes; let $t\le h$ be the number of $a_i$ equal to $1$. Then $n$ is the sum
of at most $2h$ primes and of the integer $t+2\ge2$. Every integer $u\ge2$ is a sum of at most $\frac u2$ primes: $u=2+2+\dots+2$ or $2+\dots+2+3$.
So $n$ is a sum of at most $2h+\frac{h+2}2$ primes.
\end{proof}

The proof is a template. The sieve gives an upper bound that is right up to a constant, Cauchy--Schwarz converts an upper bound for the second moment
into a lower bound for the support, and Schnirelmann's inequality converts positive density into a finite basis. Nothing in it is specific to primes
except the sieve bound. The same three steps appear in Chapter~\ref{conjectures:ch:twins}, where Erd\H{o}s used Proposition~\ref{conjectures:prop:pairs} to show that small
gaps cannot all be avoided.

\begin{remark}[the constant]
Schnirelmann's own value of $S$ was explicit and enormous. Making every step explicit was a long effort. In 1995 Olivier Ramar\'e proved that every
even number is a sum of at most six primes, so every integer greater than $1$ is a sum of at most seven. In 2012 Terence Tao proved that every odd
number greater than $1$ is a sum of at most five primes. In 2013 Harald Helfgott proved the ternary Goldbach conjecture, which Vinogradov had proved
in 1937 for all sufficiently large odd numbers: every odd number greater than $5$ is a sum of three primes. As a consequence every integer greater
than $1$ is a sum of at most four primes (Exercise~\ref{conjectures:ex:fourprimes}). The binary conjecture would make it three.

On the density side, Henry Mann proved in 1942 that $\sigma(A+B)\ge\min(1,\sigma(A)+\sigma(B))$ when $0\in A\cap B$, a conjecture of Landau and
Schnirelmann; Theorem~\ref{conjectures:thm:schnirlemma} is its weaker ancestor.
\end{remark}

\section{Why it is hard}

Two obstructions stand between Schnirelmann's theorem and Goldbach's conjecture.

\paragraph{The circle method needs three primes.} Write $e(\theta)=e^{2i\pi\theta}$ and $S(\alpha)=\sum_{p\le N}e(p\alpha)$. Since
$\int_0^1e(k\alpha)\dd\alpha$ is $1$ for $k=0$ and $0$ otherwise, expanding the powers gives
\[
r(N)=\int_0^1S(\alpha)^2e(-N\alpha)\dd\alpha,\qquad\int_0^1|S(\alpha)|^2\dd\alpha=\pi(N)
\]
(Exercise~\ref{conjectures:ex:parseval}). The method of Hardy, Littlewood and Vinogradov splits $[0,1]$ into \emph{major arcs}, short intervals around rationals
$\frac aq$ with small $q$ where $S$ can be computed from primes in progressions, and the rest, the \emph{minor arcs}, where one hopes $S$ is small. The
major arcs produce exactly the main term $\mathfrak S(N)\frac N{(\log N)^2}$. For three primes, the minor arcs contribute at most
\[
\int_{\text{minor}}|S(\alpha)|^3\dd\alpha\le\sup_{\text{minor}}|S(\alpha)|\cdot\int_0^1|S(\alpha)|^2\dd\alpha=\sup_{\text{minor}}|S(\alpha)|\cdot\pi(N),
\]
and Vinogradov's bound for $\sup|S|$, a saving of a power of $\log N$ over the trivial $\pi(N)$, makes this smaller than the main term
$\frac{N^2}{(\log N)^3}$. For two primes there is no sup to pull out: the best general bound for the minor arcs is $\int|S|^2=\pi(N)\approx\frac N{\log N}$,
which is larger than the main term $\frac N{(\log N)^2}$ by a factor $\log N$. Proving that the minor arcs do not cancel the major arcs would require
knowing that $S(\alpha)^2e(-N\alpha)$ oscillates there, and nobody knows how. What the method does give is that $r(n)>0$ for almost all even $n$: the
exceptions up to $x$ number $o(x)$, a theorem of Chudakov, van der Corput and Estermann (1937--1938).

\paragraph{Parity.} Sieves count numbers with no small prime factors, and Selberg showed that they cannot tell a number with an odd number of prime
factors from one with an even number, as explained in Chapter~\ref{conjectures:ch:twins}. A sieve that proved $r(n)>0$ would also apply to a weighted sequence
in which all the primes have been replaced by products of two primes, where the conclusion is false. The best a sieve can reach is Chen Jingrun's
theorem (1973): every large even number is $p+P_2$, with $p$ prime and $P_2$ a prime or a product of two primes.

\section{Reading a paper: Helfgott's ternary Goldbach}

Harald Helfgott's \emph{The ternary Goldbach problem} (arXiv 1501.05438, the book-length account of work first completed in May 2013) proves that every
odd number greater than $5$ is a sum of three primes. Its skeleton is the circle method of the previous section, done with complete explicitness.
\begin{itemize}
\item \emph{Minor arcs (Part I).} Bounds for $S(\alpha)$ when $\alpha$ is far from rationals with small denominators, through Vaughan's identity, which
splits $\sum\Lambda(n)e(n\alpha)$ into sums of two types, and the large sieve. Every constant is computed.
\item \emph{Major arcs (Part II).} Near $\frac aq$, $S$ is expressed through Dirichlet $L$-functions. Here Helfgott uses a rigorous computer verification,
by David Platt, of the generalised Riemann hypothesis for the relevant $L$-functions up to a finite height.
\item \emph{The integral (Part III).} The two estimates are combined, with different smoothing weights on the two sides, to show that the integral is
positive for all odd $n\ge10^{27}$; below that, a computer check (Helfgott and Platt, up to $8.8\cdot10^{30}$) finishes the job.
\end{itemize}
The idea worth stealing is the use of smoothing: replacing the sharp cutoff $p\le N$ by a smooth weight $\eta(p/N)$, chosen differently for different
parts of the argument, makes every error term smaller. Read the introduction for the history and the outline, then the chapter on the integral over the
circle, which shows how the pieces fit.

\section{Open doors}

\begin{itemize}
\item The least prime $p(n)$ with $n-p(n)$ prime grows very slowly. Compare the records of $p(n)$ with $(\log n)^2$ in the notebook. Is there a heuristic,
in the spirit of Cram\'er's model for gaps, for the size of $\max_{n\le x}p(n)$?
\item Lemoine's conjecture (1894): every odd $n\ge7$ is $p+2q$ with $p,q$ prime. Derive its Hardy--Littlewood constant and test it up to $10^7$.
\item Compute $\sigma(A)$ for $A=\{0,1\}\cup\{p+q\}$ restricted to $[1,10^6]$. Where is the infimum attained, and what $h$ does Corollary~\ref{conjectures:cor:basis}
then give?
\item The Goldbach conjecture for polynomials: every polynomial in $\mathbb F_q[t]$ of degree $\ge2$ is a sum of two irreducible polynomials? Look up
what is known, and test small cases over $\mathbb F_2$ and $\mathbb F_3$.
\end{itemize}

\section{Exercises}

\begin{exercise}[$\star$]\label{conjectures:ex:sdensity}
Compute the Schnirelmann density of the odd numbers and of the squares. Show that $\sigma(A)=1$ if and only if $A$ contains every positive integer, and
that $\sigma(A)=0$ if $1\notin A$.
\end{exercise}

\begin{exercise}[$\star\star$]\label{conjectures:ex:basisorder}
Let $0\in A$ with $\sigma(A)=\alpha>0$. Show that every nonnegative integer is a sum of at most $2\bigl\lceil\frac{\log2}{-\log(1-\alpha)}\bigr\rceil$ elements
of $A$ (for $\alpha=1$ read the fraction as $0$ and the bound as $1$). What does this give for $\alpha=\frac12$?
\end{exercise}

\begin{exercise}[$\star$]\label{conjectures:ex:goldbachforms}
Show that Euler's form of the conjecture implies that every odd number greater than $5$ is a sum of three primes, and that it is equivalent to
Goldbach's marginal form (in modern terms): every integer greater than $5$ is a sum of three primes.
\end{exercise}

\begin{exercise}[$\star$]\label{conjectures:ex:oddr}
Show that for odd $n$, $r(n)=2$ if $n-2$ is prime and $r(n)=0$ otherwise.
\end{exercise}

\begin{exercise}[$\star$]\label{conjectures:ex:smallgoldbach}
Check that every even $n$ with $4\le n\le100$ is $p+q$ with $p\le19$, and that $98$ needs $p=19$.
\end{exercise}

\begin{exercise}[$\star$]\label{conjectures:ex:bands}
Compute $\mathfrak S(n)/\mathfrak S(2^k)$ for $n=2^a3^b$, $n=2^a3^b5^c$ and $n=2^a3^b5^c7^d$ with all exponents positive. Explain the bands of
Figure~\ref{conjectures:fig:comet}, and say which even numbers form the lowest edge of the comet.
\end{exercise}

\begin{exercise}[$\star\star$]\label{conjectures:ex:singavg}
Show that the average of $\mathfrak S(n)$ over the even $n\le x$ tends to $2$. Check that this is consistent with $R(x)=\sum_{m\le x}r(m)\sim\frac{x^2}{2(\log x)^2}$.
\end{exercise}

\begin{exercise}[$\star$]\label{conjectures:ex:parseval}
Prove the two identities $r(N)=\int_0^1S(\alpha)^2e(-N\alpha)\dd\alpha$ and $\int_0^1|S(\alpha)|^2\dd\alpha=\pi(N)$. Find the corresponding formula for the
number of representations of $N$ as a sum of three primes.
\end{exercise}

\begin{exercise}[$\star$]\label{conjectures:ex:fourprimes}
Using Helfgott's theorem, show that every integer $n\ge2$ is a sum of at most four primes, and that four cannot be replaced by two.
\end{exercise}

\begin{exercise}[$\star\star$]\label{conjectures:ex:sumbounds}
Show that $\pi(\frac x2)^2\le R(x)\le\pi(x)^2$. Deduce from the prime number theorem that $R(x)\asymp\frac{x^2}{(\log x)^2}$, and then that $r(m)$ is at least
$\frac{c\,x}{(\log x)^2}$ for a positive proportion of the even $m\le x$, using only Proposition~\ref{conjectures:prop:goldbachsieve} and Lemma~\ref{conjectures:lem:phifour}.
\end{exercise}

\section{Hints}

\begin{hint}{conjectures:ex:sdensity}
For the odd numbers look at even $n$; for the squares, $A(n)\le\sqrt n$.
\end{hint}
\begin{hint}{conjectures:ex:basisorder}
Follow the proof of Corollary~\ref{conjectures:cor:basis} with the least possible $k$.
\end{hint}
\begin{hint}{conjectures:ex:goldbachforms}
Subtract $2$ or $3$. For the converse, a sum of three primes that is even contains the prime $2$.
\end{hint}
\begin{hint}{conjectures:ex:oddr}
A sum of two odd primes is even.
\end{hint}
\begin{hint}{conjectures:ex:smallgoldbach}
Only $98$ is hard: try $p=2,3,5,\dots$ in turn.
\end{hint}
\begin{hint}{conjectures:ex:bands}
Only odd prime divisors of $n$ change $\mathfrak S(n)$, through the factor $\frac{p-1}{p-2}$.
\end{hint}
\begin{hint}{conjectures:ex:singavg}
Expand $\prod_{p\mid n,p>2}\frac{p-1}{p-2}=\sum_{d\mid n,\,d\text{ odd}}\frac{\mu^2(d)}{\psi(d)}$, where $\psi$ is multiplicative with $\psi(p)=p-2$.
\end{hint}
\begin{hint}{conjectures:ex:parseval}
Expand $S(\alpha)^2$ as a double sum over $p,q$.
\end{hint}
\begin{hint}{conjectures:ex:fourprimes}
For $n\ge9$ use $n-3$; check the small cases by hand. For the second part, look at $n=27$.
\end{hint}
\begin{hint}{conjectures:ex:sumbounds}
For the last part, rerun the Cauchy--Schwarz step of Theorem~\ref{conjectures:thm:schnirelmann} on the $m$ with $r(m)\ge\frac{R(x)}{2x}$.
\end{hint}

\section{Solutions}

\begin{solution}{conjectures:ex:sdensity}
\idea Schnirelmann density is an infimum over all $n$, so one bad $n$ decides it.
\writeup For the odd numbers $A(n)=\lceil\frac n2\rceil\ge\frac n2$, with equality for even $n$, so $\sigma=\frac12$. For the squares $A(n)=\lfloor\sqrt n\rfloor$, so
$\frac{A(n)}n\le\frac1{\sqrt n}\to0$ and $\sigma=0$. If $\sigma(A)=1$ then $A(n)\ge n$, so $A(n)=n$ for every $n$, which means $[1,n]\subset A$ for every $n$; the
converse is clear. If $1\notin A$ then $\frac{A(1)}1=0$.
\end{solution}

\begin{solution}{conjectures:ex:basisorder}
\idea Make the choice of $k$ in Corollary~\ref{conjectures:cor:basis} explicit.
\writeup We need $(1-\alpha)^k\le\frac12$, that is $k\log(1-\alpha)\le-\log2$, that is $k\ge\frac{\log2}{-\log(1-\alpha)}$ when $\alpha<1$. The least such $k$ is
the ceiling, and $h=2k$ elements suffice (since $0\in A$, sums of fewer elements are sums of exactly $h$). For $\alpha=1$, $A$ contains every positive integer
and one element suffices. For $\alpha=\frac12$, $k=1$, and Lemma~\ref{conjectures:lem:cover} applies directly: $A+A$ is everything. The odd numbers with $0$ show this
is sharp: $\sigma=\frac12$, and $2$ is not in $A$, so one element is not enough.
\end{solution}

\begin{solution}{conjectures:ex:goldbachforms}
\idea Parity decides how many $2$'s a representation must contain.
\writeup Assume every even $n\ge4$ is $p+q$. If $n>5$ is odd, $n-3\ge4$ is even, so $n=3+p+q$; if $n>5$ is even, $n-2\ge4$ is even, so $n=2+p+q$. Conversely,
assume every integer greater than $5$ is a sum of three primes, and let $n\ge4$ be even. If $n=4$, $n=2+2$. Otherwise $n+2>5$, so $n+2=p+q+r$. The sum is
even, so the number of odd primes among $p,q,r$ is even: either none, and $n+2=6$, $n=4$, already done; or two, and the third is $2$, so $n=p+q$ for the two
odd ones.
\end{solution}

\begin{solution}{conjectures:ex:oddr}
\idea One of the two primes must be even.
\writeup If $p+q=n$ is odd, $p$ and $q$ have different parities, so one of them is $2$. The ordered pairs are $(2,n-2)$ and $(n-2,2)$, which count when $n-2$
is prime; they are distinct because $n-2\ne2$ for odd $n$. So $r(n)=2$ if $n-2$ is prime and $0$ otherwise.
\end{solution}

\begin{solution}{conjectures:ex:smallgoldbach}
\idea Search upwards from the smallest primes.
\writeup For each even $n\le100$ the least $p$ with $n-p$ prime is quickly found: it is $2$ for $n=4$, $3$ for most $n$, and the larger values are
$p(12)=5$, $p(30)=7$, $p(98)=19$ (the records in this range, from the table computed in the notebook). For $98$: $98-2=96$, $98-3=95=5\cdot19$,
$98-5=93=3\cdot31$, $98-7=91=7\cdot13$, $98-11=87=3\cdot29$, $98-13=85=5\cdot17$, $98-17=81=3^4$ are composite, and $98-19=79$ is prime.
\end{solution}

\begin{solution}{conjectures:ex:bands}
\idea The singular series only sees the set of odd primes dividing $n$.
\writeup Since $\mathfrak S(n)=2C_2\prod_{p\mid n,p>2}\frac{p-1}{p-2}$, the ratios are $2$ for $2^a3^b$, $2\cdot\frac43=\frac83$ for $2^a3^b5^c$, and
$\frac83\cdot\frac65=\frac{16}5$ for $2^a3^b5^c7^d$. The factor for $p$ is $1+\frac1{p-2}$, so the prime $3$ doubles the prediction while large primes barely
change it. That splits the comet into a band for $3\nmid n$ and a band twice as high for $3\mid n$, each of which is subdivided more finely by $5$, $7$,
and so on. The lowest edge consists of the $n$ with no small odd prime factors, the extreme case being the powers of $2$, where $\mathfrak S=2C_2$.
\end{solution}

\begin{solution}{conjectures:ex:singavg}
\idea Expand the product as a divisor sum and swap, as for $\frac n{\varphi(n)}$ in Chapter~\ref{conjectures:ch:twins}.
\writeup Write $n=2m$ with $m\le X=\frac x2$. Since $\frac{p-1}{p-2}=1+\frac1{p-2}$, we have $\prod_{p\mid n,p>2}\frac{p-1}{p-2}=\sum_{d\mid m,\,d\text{ odd}}\frac{\mu^2(d)}{\psi(d)}$,
using that an odd $d$ divides $2m$ exactly when it divides $m$. So
\[
\sum_{m\le X}\prod_{\substack{p\mid2m\\p>2}}\frac{p-1}{p-2}=\sum_{\substack{d\le X\\d\text{ odd}}}\frac{\mu^2(d)}{\psi(d)}\Bigl\lfloor\frac Xd\Bigr\rfloor=X\sum_{d\text{ odd}}\frac{\mu^2(d)}{d\psi(d)}+O\Bigl(X\sum_{d>X}\frac{\mu^2(d)}{d\psi(d)}+\sum_{d\le X}\frac{\mu^2(d)}{\psi(d)}\Bigr).
\]
For odd $p$, $p-2\ge\frac p3$, so $\frac{\mu^2(d)}{\psi(d)}\le\frac{3^{\omega(d)}}d$ and $\frac1{d\psi(d)}\le\frac{3^{\omega(d)}}{d^2}$. Hence
$\sum_{d\le X}\frac{\mu^2(d)}{\psi(d)}\le\prod_{p\le X}(1+\frac3p)\le\exp\bigl(3\sum_{p\le X}\frac1p\bigr)=O\bigl((\log X)^3\bigr)$ by Mertens (Tome~5,
Theorem~1.5). For the tail, $\frac1{d^2}\le\frac1{\sqrt X\,d^{3/2}}$ when $d>X$, so it is at most
$X^{-1/2}\sum_d\frac{\mu^2(d)3^{\omega(d)}}{d^{3/2}}=X^{-1/2}\prod_p\bigl(1+\frac3{p^{3/2}}\bigr)=O(X^{-1/2})$. Both errors are $o(X)$. The series is
$\prod_{p>2}\bigl(1+\frac1{p(p-2)}\bigr)=\prod_{p>2}\frac{(p-1)^2}{p(p-2)}=\frac1{C_2}$. So the average of $\mathfrak S(n)$ over even $n\le x$ tends to
$2C_2\cdot\frac1{C_2}=2$.

Consistency: if $r(m)\approx\mathfrak S(m)\frac m{(\log m)^2}$ for even $m$, then $R(x)\approx\frac2{(\log x)^2}\sum_{m\le x,\,m\text{ even}}m\approx\frac2{(\log x)^2}\cdot\frac{x^2}4$,
which is $\frac{x^2}{2(\log x)^2}$. And indeed $R(x)$ counts pairs of primes in the triangle $p+q\le x$, whose expected number is about
$\iint_{s+t\le x}\frac{\dd s\,\dd t}{\log s\log t}\sim\frac{x^2}{2(\log x)^2}$.
\end{solution}

\begin{solution}{conjectures:ex:parseval}
\idea Orthogonality: $\int_0^1e(k\alpha)\dd\alpha=[k=0]$.
\writeup $S(\alpha)^2e(-N\alpha)=\sum_{p,q\le N}e((p+q-N)\alpha)$; integrating term by term keeps the pairs with $p+q=N$, which is $r(N)$ since such pairs have
$p,q\le N$ anyway. Similarly $|S(\alpha)|^2=S(\alpha)\overline{S(\alpha)}=\sum_{p,q\le N}e((p-q)\alpha)$, and the integral keeps the pairs $p=q$: there are
$\pi(N)$ of them. For three primes, $\int_0^1S(\alpha)^3e(-N\alpha)\dd\alpha$ counts the ordered triples of primes with $p+q+r=N$.
\end{solution}

\begin{solution}{conjectures:ex:fourprimes}
\idea Subtract a small prime to land in Helfgott's range.
\writeup The integers $2,3$ are primes, and $4=2+2$, $5=2+3$, $6=3+3$, $7=2+5$, $8=3+5$. For $n\ge9$, $n-3\ge6$. If $n$ is even, $n-3$ is odd and greater than
$5$, so it is a sum of three primes, and $n$ is a sum of four. If $n$ is odd, $n$ itself is a sum of three primes. Two primes are not enough: $27$ is odd and
$27-2=25$ is not prime, so $r(27)=0$ by Exercise~\ref{conjectures:ex:oddr}, and $27$ is not prime.
\end{solution}

\begin{solution}{conjectures:ex:sumbounds}
\idea Pairs in a square inside the triangle, the triangle inside a bigger square; then Cauchy--Schwarz on the large values.
\writeup Every pair with $p,q\le\frac x2$ has $p+q\le x$, and every pair with $p+q\le x$ has $p,q\le x$; this gives the two bounds, and the prime number theorem
makes both of order $\frac{x^2}{(\log x)^2}$. Let $\mathcal M$ be the set of $m\le x$ with $r(m)\ge\frac{R(x)}{2x}$. The $m\notin\mathcal M$ contribute at most
$x\cdot\frac{R(x)}{2x}=\frac{R(x)}2$ to $R(x)$, so $\sum_{m\in\mathcal M}r(m)\ge\frac{R(x)}2$, and by Cauchy--Schwarz and the bound
$\sum_{m\le x}r(m)^2\le C_4\frac{x^3}{(\log x)^4}$ from the proof of Theorem~\ref{conjectures:thm:schnirelmann},
\[
|\mathcal M|\ge\frac{R(x)^2/4}{\sum_{m\le x}r(m)^2}\ge c'x.
\]
Every $m\in\mathcal M$ has $r(m)\ge\frac{R(x)}{2x}\ge\frac{c\,x}{(\log x)^2}$, and $m$ is even unless $r(m)\le2$, which is excluded for large $x$.
\end{solution}

\section*{Sources and further reading}

\begin{itemize}
\item C.~Goldbach to L.~Euler, 7 June 1742, and Euler's reply of 30 June 1742, in P.-H.~Fuss (ed.), \emph{Correspondance math\'ematique et physique de
quelques c\'el\`ebres g\'eom\`etres du XVIII\`eme si\`ecle}, vol.~1, St Petersburg, 1843 (the epigraph translates Euler's ``ein jeder numerus par eine summa
duorum primorum sey, halte ich f\"ur ein ganz gewisses theorema'', with the clause that follows it, in our words).
\item G.~H.~Hardy and J.~E.~Littlewood, \emph{Some problems of `Partitio numerorum' III: on the expression of a number as a sum of primes}, Acta
Mathematica 44 (1923), 1--70 (Conjecture~A, the asymptotic formula for $r(n)$).
\item T.~Oliveira e Silva, S.~Herzog and S.~Pardi, \emph{Empirical verification of the even Goldbach conjecture and computation of prime gaps up to
$4\cdot10^{18}$}, Mathematics of Computation 83 (2014), 2033--2060; the project page of T.~Oliveira e Silva for the dates and the record $p(n)=9781$.
J.~Richstein, \emph{Verifying the Goldbach conjecture up to $4\cdot10^{14}$}, Mathematics of Computation 70 (2001), 1745--1749, and M.~K.~Sinisalo,
\emph{Checking the Goldbach conjecture up to $4\cdot10^{11}$}, Mathematics of Computation 61 (1993), 931--934.
\item L.~Schnirelmann, \emph{\"Uber additive Eigenschaften von Zahlen}, Mathematische Annalen 107 (1933), 649--690 (the German account of his 1930 work).
The MacTutor biography of Lev Schnirelmann (University of St Andrews) for his life.
\item H.~B.~Mann, \emph{A proof of the fundamental theorem on the density of sums of sets of positive integers}, Annals of Mathematics 43 (1942),
523--527.
\item O.~Ramar\'e, \emph{On \v{S}nirel'man's constant}, Annali della Scuola Normale Superiore di Pisa, Classe di Scienze (4) 22 (1995), 645--706.
T.~Tao, \emph{Every odd number greater than $1$ is the sum of at most five primes}, Mathematics of Computation 83 (2014), 997--1038.
\item I.~M.~Vinogradov, \emph{Representation of an odd number as a sum of three primes}, Doklady Akademii Nauk SSSR 15 (1937), 291--294.
H.~A.~Helfgott, \emph{The ternary Goldbach problem}, arXiv:1501.05438 (2015), the paper of this chapter's reading guide.
\item J.~R.~Chen, \emph{On the representation of a larger even integer as the sum of a prime and the product of at most two primes}, Scientia Sinica 16
(1973), 157--176.
\item ``Gold for Goldbach'', \emph{Plus Magazine}, June 2000 (the Faber and Faber prize); A.~Doxiadis, \emph{Uncle Petros and Goldbach's Conjecture},
Faber and Faber, 2000, the novel it promoted.
\item M.~B.~Nathanson, \emph{Additive Number Theory: The Classical Bases}, Springer, 1996, for Schnirelmann's theorem and its refinements.
\end{itemize}
